\section{总结与展望}

本研究围绕课程知识的命名实体识别与关系抽取，研究了大语言模型在计算机网络课程知识图谱自动化构建中的应用，完成的工作包括以下几个方面：

（1）在课程语料构建与预处理方面，实现了一套面向计算机网络课程的多源语料自动化预处理流程。语料覆盖教材（谢希仁《计算机网络》第8版\cite{xiexiren2021}、Kurose\&Ross《计算机网络：自顶向下方法》第8版、《TCP/IP详解》等）、课堂讲义（9章课件及10次授课讲稿）、实验指导书、历年试题以及50余份RFC标准文档，共205份PDF文档。预处理阶段使用PyMuPDF渲染PDF页面，LibreOffice无头模式处理PPT/PPTX格式文档，Qwen3-VL-8B多模态视觉语言模型执行OCR文字提取，最终输出结构化Markdown文本。该语料库覆盖多种来源与格式，为后续知识抽取阶段提供了基础。

（2）在命名实体识别方面，本课题为计算机网络领域定义了四类实体：协议名称（protocol，如TCP、HTTP、OSPF）、概念术语（concept，如拥塞控制、子网掩码、分层模型）、机制算法（mechanism，如三次握手、慢启动、CSMA/CD）和性能指标（metric，如吞吐量、RTT、丢包率）。考虑到课程领域缺乏标注数据，研究了零样本（zero-shot）和少样本（few-shot）两种提示策略。零样本模式通过详细的实体类型定义与标注指南引导模型行为；少样本模式在此基础上增加了四条覆盖全部关系类型的标注示例以及负面示例，改进模型对课程术语边界的识别。在启用LLM复核（\texttt{fs-recheck}）的最终配置下抽取并合并得到32,760个实体，频次跨度从1次到2,117次（平均3.32次），高频实体如TCP（2,117次）、UDP（1,069次）、HTTP（701次）、IP（640次）等，分布与计算机网络课程的知识结构大体一致。

（3）在关系抽取方面，本课题定义了四类关系类型：包含关系（contains），表达概念或协议间的结构组成；依赖关系（depends\_on），表示功能实现上的前置条件；对比关系（compared\_with），描述同类实体间的异同；应用关系（applied\_to），说明机制算法在具体协议或场景中的部署方式。抽取架构上，设计了Propose--Review“提议—审核”两阶段流程：第一阶段由大语言模型按Schema生成候选三元组；第二阶段通过独立审核提示词进行质量门控，逐条检查证据是否为输入文本的连续原文子串、头尾实体是否均出现在证据中、关系类型是否合法。该机制要求每条三元组给出可回溯的原文证据，减少了模型生成不实内容的情况。知识融合阶段，通过别名映射表（如“传输控制协议”→“TCP”）以及基于FAISS索引的嵌入语义规范化（结合UnionFind算法）实现跨文档实体消歧。最终知识图谱（\texttt{merged-fs-recheck}）包含32,760个节点和17,627条边，连通子图（13,018个节点）平均度数约2.71。在500条独立金标上采用“严格／软匹配／槽位／类型级”四档匹配口径评估，软匹配三元组F1为15\%~19\%、槽位F1为37\%~42\%、类型级F1可达74\%~82\%。从严格档到类型级档的阶梯说明系统失分主要源于实体表面形式差异，而非关系类型识别错误。

（4）在知识抽取管线方面，本课题构建了六阶段端到端自动化流程：文档渲染（PDF$\rightarrow$PNG）$\rightarrow$OCR文字提取（PNG$\rightarrow$Markdown）$\rightarrow$文本分块（Markdown$\rightarrow$JSON Chunks）$\rightarrow$三元组抽取（Chunks$\rightarrow$Raw Triples）$\rightarrow$知识融合（Raw Triples$\rightarrow$Merged Triples）$\rightarrow$图数据库导入（Merged Triples$\rightarrow$Neo4j）。流水线采用工厂模式统一管理LLM调用，通过滑动窗口速率限制和指数退避重试保障API调用稳定性。各阶段经中间文件解耦，支持断点续跑与独立调试。图存储层面，融合后的三元组通过Cypher MERGE语句幂等导入Neo4j，配合 \texttt{name} 唯一约束与 \texttt{frequency} 索引为下游应用提供基于实体类型、关系类型与频次的高效检索能力。系统验证了大语言模型在课程知识建模中的可行性。

